\documentclass[../../../include/open-logic-section]{subfiles}

\begin{document}

\olfileid{sth}{choice}{countablechoice}
\olsection{లెక్కించదగిన ఎంపిక}

ఎంపిక/సుక్రమపరచడాన్ని ఏమాత్రం వాడకుండానే ఈ విషయాన్ని సులభంగా నిరూపించవచ్చు:

\begin{lem}[$\Zminus$లో]
ప్రతి పరిమిత సమితికీ ఎంపిక ప్రమేయం ఉంది.
\end{lem}

\begin{proof}
$a = \{b_1, \ldots, b_n\}$ అనుకుందాం; ఇక్కడ సభ్యులను పునరావృత్తి లేకుండా జాబితా చేసుకున్నాం. సరళత కోసం ప్రతి $b_i
\neq \emptyset$ అనుకుందాం. అప్పుడు $c_1 \in b_1, \ldots, c_n \in b_n$ అయ్యే $c_1, \ldots, c_n$ వస్తువులు ఉన్నాయి. ఇప్పుడు \olref[z][pairs]{prop:pairsconsequences} ప్రకారం $\{\langle b_1,
c_1\rangle , \ldots, \langle b_n, c_n\rangle\}$ అనే సమితి ఉంది; అదే~$a$ కోసం ఎంపిక ప్రమేయం.
\sourcecorrection{OLTESTCHOICECOUNT-001}{మూలం $a$ సభ్యుల జాబితాలో పునరావృత్తి లేదని చెప్పలేదు. ఒకే $b_i$ పలుసార్లు ఉండి వేర్వేరు $c_i$లను ఎంచుకుంటే చూపిన సంబంధం ప్రమేయం కాకపోవచ్చు; వేర్వేరు సభ్యుల జాబితా అని స్పష్టం చేశాం.}
\end{proof}

కానీ అనంత సమితులను చూడగానే విషయాలు క్లిష్టమవుతాయి. ఉదాహరణకు, పై ఫలితానికి ఈ ``కనిష్ఠ'' విస్తరణను చూడండి:

\begin{defish}
\emph{లెక్కించదగిన ఎంపిక.} ప్రతి \emph{లెక్కించదగిన} సమితికీ ఎంపిక ప్రమేయం ఉంది.
\end{defish}

ఇది ఎంపికకు ప్రత్యేక సందర్భం. ఈ సూత్రాన్ని వాడుతున్నామనే స్పష్టమైన అవగాహన లేకుండానే తరచుగా దాన్ని ఉపయోగించారని తేలుతుంది. రెండు మంచి ఉదాహరణలు ఇవి.\footnote{\citet[\S9.4]{Potter2004}, లూకా ఇంకుర్వాటీ నుంచి.}

\begin{ex}
ఒక సహజ ఆలోచన ఇది: ఏ సమితి $A$కైనా $\cardle{\omega}{A}$ లేదా ఏదో $n \in \omega$కు $\cardeq{A}{n}$. ప్రతి సమితీ పరిమితమో అనంతమో అన్న సహజ భావనను చెప్పే ఒక మార్గం ఇది. కాంటర్, మరెందరో గణితవేత్తలు నిరూపణ లేకుండానే ఈ వాదన చేశారు. మనం జాగ్రత్తగా \olref[cardinals][classing]{generalinfinitycharacter}లో నిరూపించాం. కానీ అక్కడ ప్రతి సమితి $A$ను సుక్రమపరచవచ్చనీ, అందువల్ల $\card{A}$ తప్పక ఉంటుందనీ అనుకుని $\ZFC$లో పనిచేశాం. అంటే ఎంపికను స్పష్టంగానే అనుకున్నాం.

నిజానికి \citet{Dedekind1888} ఈ వాదనకు తన నిరూపణను ఇలా ఇచ్చారు:

\begin{thm}[$\Zminus + \text{Countable Choice}$లో]
ఏ $A$కైనా $\cardle{\omega}{A}$ లేదా ఏదో $n \in \omega$కు $\cardeq{A}{n}$.
\end{thm}

\begin{proof}
అన్ని $n \in \omega$లకు $\cardneq{A}{n}$ అనుకుందాం. అప్పుడు ప్రతి $n < \omega$కు ఖచ్చితంగా $\cardexpo{2}{n}$ మూలకాలున్న ఉపసమితి $A_n \subseteq A$ ఉంటుంది. ఈ $A_0, A_1, A_2,
\ldots$ శ్రేణిని వాడి ప్రతి $n$కు ఇలా నిర్వచిద్దాం:
\[
	B_n = A_n \setminus \bigcup_{i < n} A_i.
\]
ఇప్పుడు గమనించండి:
\begin{align*}
	\card{\bigcup_{i < n}A_i} 
	&\leq \card{A_0} + \card{A_1} + \ldots + \card{A_{n-1}}\\
	&=1 + 2 + \ldots + 2^{n-1}\\
	& = 2^n - 1\\
	& < 2^n = \card{A_n}
\end{align*}
కాబట్టి ప్రతి $B_n$లో కనీసం ఒక మూలకం $c_n$ ఉంది. ఇంకా $B_n$లు జతలవారీ విచ్ఛిన్నాలు; కనుక $c_n = c_m$ అయితే $n = m$. అలాగే ప్రతి $c_n
\in A$. కాబట్టి $f(n) = c_n$ అనే ప్రమేయం $\omega \to A$ ఇంజెక్షన్.
\sourcecorrection{OLTESTCHOICECOUNT-002}{మూల అసమానత ఎడమవైపు $\bigcup_{i<n}A_n$ అని, మారే సూచిక $i$ బదులు స్థిర $A_n$ను వ్రాసింది; అది తరువాతి $<\card{A_n}$కి విరుద్ధం. పూర్వ సమితుల సమ్మేళనం $\bigcup_{i<n}A_i$గా సరిచేశాం.}
\end{proof}
\noindent
డెడెకిండ్ తాను లెక్కించదగిన ఎంపికను వాడానని చెప్పలేదు. కానీ \emph{మీరు} దాని వినియోగాన్ని గుర్తించారా? మళ్లీ చూడండి. (నిజంగా: \emph{మళ్లీ చూడండి}.)

ఈ నిరూపణ లెక్కించదగిన ఎంపికను రెండుసార్లు వాడింది. $A_0$, $A_1$, $A_2$, \dots\@ సమితుల శ్రేణిని పొందడానికి ఒకసారి; ప్రతి~$B_n$ నుంచి $c_n$ మూలకాన్ని ఎంచుకోవడానికి మరోసారి. అంతేకాదు, ఈ ఎంపిక వినియోగాన్ని తొలగించలేం. ఎంపిక రూపం ఏదీ లేకపోతే ఈ ఫలితం విఫలమవుతుందని \citet[p.~138]{Cohen1966} నిరూపించారు. అంటే $\omega$తో \emph{పోల్చలేని} సమితులు ఉండటం $\ZF$తో సుసంగతం.
\end{ex}

\begin{ex}
లెక్కించదగిన సమితుల లెక్కించదగిన సమ్మేళనం లెక్కించదగినదే అని \citeyear{Cantor1878}లో కాంటర్ చెప్పారు. ఆయన నిరూపణ ఇవ్వలేదు; బహుశా అది స్పష్టమే అనుకున్నారు. మనం జాగ్రత్తగా \olref[card-arithmetic][simp]{kappaunionkappasize}లో దీని మరింత సాధారణ రూపాన్ని నిరూపించాం. కానీ మన నిరూపణ ఎంపికను స్పష్టంగా అనుకుంది. అంత సాధారణం కాని ఫలితానికీ లెక్కించదగిన ఎంపిక అవసరం.

\begin{thm}[$\Zminus + \text{Countable Choice}$లో]
ప్రతి $n \in \omega$కు $A_n$ లెక్కించదగినదైతే, $\bigcup_{n <
\omega} A_n$ కూడా లెక్కించదగినది.
\end{thm}

\begin{proof}
సాధారణతకు భంగం లేకుండా ప్రతి $A_n \neq \emptyset$ అనుకుందాం. అప్పుడు ప్రతి $n \in \omega$కు $f_n \colon \omega
\to A_n$ అనే \tetoken{సర్జెక్షన్}{!!a{surjection}} ఉంటుంది. $f(m, n) = f_n(m)$గా $f \colon \omega \times \omega \to \bigcup_{n <
\omega} A_n$ను నిర్వచిద్దాం. $\omega
\times \omega$ లెక్కించదగినది (\olref[sfr][siz][zigzag]{natsquaredenumerable}), $f$ \tetoken{సర్జెక్షన్}{!!a{surjection}} కాబట్టి ఫలితం వస్తుంది.
\end{proof}
\noindent
లెక్కించదగిన ఎంపిక ఎక్కడ వాడామో గుర్తించారా? $f_0$, $f_1$, $f_2$, \dots\@ ప్రమేయాల శ్రేణిని ఎంచుకోవడానికి.\footnote{\olref[card-arithmetic][simp]{kappaunionkappasize}లో ``ప్రతి $\beta \in \cardfont{a}$కు $f_\beta$ అనే \tetoken{ఇంజెక్షన్‌ను}{!!a{injection}} ఎంచుకోండి'' అన్నప్పుడూ ఇదే విధమైన ఎంపికను వాడాం.} మళ్లీ, ఎంపిక సూత్రం ఏదీ లేకపోతే ఫలితం విఫలమవుతుంది. ముఖ్యంగా లెక్కించదగిన సమితుల లెక్కించదగిన సమ్మేళనానికి~$\beth_1$ కార్డినాలిటీ ఉండటం $\ZF$తో సుసంగతమని \citet{FefermanLevy1963} నిరూపించారు. ఇదే విషయాన్ని రస్సెల్ మరింత వినోదాత్మకంగా ఇలా చెప్పారు:
\begin{quote}
  ఒక కోటీశ్వరుడు బూట్ల జత కొన్న ప్రతిసారీ మేజోళ్ల జత కూడా కొనేవాడు; మరే సందర్భంలోనూ కొనేవాడు కాదు. రెండింటినీ కొనడమంటే అతనికి ఎంత ఇష్టమంటే చివరికి అతని దగ్గర $\aleph_0$ బూట్ల జతలు, $\aleph_0$ మేజోళ్ల జతలు వచ్చాయి\dots\@ బూట్లలో కుడి, ఎడమను వేరు చేయవచ్చు. కాబట్టి ప్రతి జత నుంచి ఒకదాన్ని, ఉదాహరణకు అన్ని కుడి బూట్లను లేదా అన్ని ఎడమ బూట్లను, ఎంచుకోవచ్చు. కానీ మేజోళ్లకు అలాంటి ఎంపిక నియమం తోచదు; గుణిత స్వయంసిద్ధాన్ని [అంటే ప్రభావంలో ఎంపికను] అనుకోకపోతే, ప్రతి జత నుంచి ఒక్కో మేజోళ్లు గల తరగతి ఉందని నిశ్చయంగా చెప్పలేం.
  \citep[p.~126]{Russell1919}
\end{quote}
సంక్షిప్తంగా, లెక్కించదగిన సంఖ్యలో మేజోళ్ల జతలు ఉంటే, లెక్కించదగిన సంఖ్యలోనే మేజోళ్లు ఉంటాయని నిరూపించడానికి ఎంపిక యొక్క ఏదో రూపం అవసరం. నిజానికి లెక్కించదగిన ఎంపిక (లేదా దానికి సమానమైనది) లేకపోతే, లెక్కించదగిన సమితుల లెక్కించదగిన సమ్మేళనం లెక్కించదగినది కాకపోవచ్చు.
\end{ex}

దీనిలోని పాఠం: లెక్కించదగిన ఎంపికను వాడుతున్నామని చాలా మందికి తెలియకుండానే దాన్ని పదేపదే వాడారు. తాత్త్విక ప్రశ్న: లెక్కించదగిన ఎంపికను ఎలా \emph{సమర్థించవచ్చు}?

సహజమైన సమర్థన ప్రయత్నం ఒక అతికార్యాన్ని ఆశ్రయించవచ్చు. మొదటి ఎంపికను ఒక నిమిషంలో $\nicefrac{1}{2}$ భాగంలో, రెండవ ఎంపికను $\nicefrac{1}{4}$ భాగంలో, \dots, $n$వ ఎంపికను $\nicefrac{1}{2^n}$ భాగంలో, \dots\@ చేద్దాం. అప్పుడు $1$~నిమిషంలో $\omega$-శ్రేణి ఎంపికలను చేసి, ఎంపిక ప్రమేయాన్ని నిర్వచించినట్టవుతుంది.

కానీ ఇలాంటి ఆలోచనా ప్రయోగం నిజంగా ఏం చెబుతుంది? మొదట, ``ఎంచుకోవడం'' అనే భావనను అక్షరార్థంలో తీసుకుంటుంది. రెండవది, గణితాన్ని అధిభౌతిక సాధ్యతతో ముడిపెడుతుంది.

ముఖ్యంగా, ఇది \emph{కేవలం} లెక్కించదగిన ఎంపికకు మించి ఎంపికను \emph{పూర్తిగా} సమర్థించదు. \emph{ప్రతి} సమితికీ ఎంపిక ప్రమేయం కావాలంటే, ``ఏ క్రమసంఖ్య పొడవుకైనా అతికార్యం'' చేయగలగాలి. స్పష్టంగా చెప్పాలంటే, ఆ ఆలోచన హాస్యాస్పదం.

\end{document}
